\scanpage{114}
\mspage{113}
\uline{At M 4} Si $\Phi,\Psi\in\unsure{\widetilde{\mathcal{F}}|\widetilde{\mathscr{F}}}$, alors $\Phi\cup\Psi\in\widetilde{\mathcal{F}}$ ssi
\[
\begin{aligned}
&\forall X\in\Phi,\ Y\in\Psi,\ \deleted{\unsure{\text{on a}}\ X,Y\notin\Phi\cap\Psi,\ \text{on a }X\mathrel{\unsure{|\circ|\,|\,|o|}}Y.}\\
&\inserted{\text{on a }X\mathrel{\unsure{\gtrless|\lessgtr}}Y\ \text{i.e.}}\\
&\inserted{\forall X'\in\widetilde{X}\setminus\widetilde{X}\cap\widetilde{Y},\ Y'\in\widetilde{Y}\setminus\widetilde{X}\cap\widetilde{Y},\ \text{on a }\deleted{\unsure{X\mathrel{|\circ|}Y}}X'\mathrel{|\circ|}Y'.}
\end{aligned}
\]

L'atelier est dit \deleted{fini} ``à figures finies'' si les maquettes de ses figures sont finies. Cela revient \unsure{à|:} \unsure{se} donner $\deleted{\unsure{M}}$, le magasin correspondant, on trouve alors
\[
\widetilde{\mathcal{F}}=\{\deleted{\unsure{A}}\Phi\in\unsure{P|\mathcal{P}}_{f}(M)\mid X,Y\in\deleted{\unsure{A}}\Phi\Longrightarrow X\mathrel{\gtrless}Y\},
\]
mais \deleted{\unsure{est}} cela satisfait aux conditions d'un $M$-atelier ssi on a

\deleted{\unsure{At M}}\par
\uline{Mag fin} $\forall X\in M$, $\widetilde{X}$ est fini.

(Axiome des \uline{magasins} à \deleted{\unsure{multistrates finies}}\inserted{\uline{maquettes finies}}).

Au cours du travail, \inserted{dans cette petite \unsure{algèbre},} c'est finalement \uline{l'\unsure{étalier|atelier}} qui est au centre de l'attention constamment, alors que l'ensemble $\mathcal{F}$ de toutes les figures joue un rôle \unsure{strictement} accessoire.
